\chapter{The EPL Language}
\label{chap:language}

This chapter describes the \epl{} language from the programmer's point of view.
The implementation is intentionally small, so the language can be explained
without hiding large subsystems. Even so, \epl{} includes enough features to
compile useful examples: functions, external declarations, local variables,
expression blocks, conditionals, loops, structs, arrays, raw pointers, casts, and
compile-time evaluation.

\section{Design Goals}

The first design goal is simplicity. \epl{} is a single-file language. There are
no modules, imports, packages, namespaces, generics, traits, classes, closures,
or overloads. Function names are global and are emitted as written. Struct names
live in a separate global type namespace. This keeps name resolution and code
generation understandable.

The second design goal is explicitness. Numeric conversions are not implicit.
Pointer operations are written explicitly. External functions are declared
explicitly. A function's parameter types are always written in the source. The
compiler performs local type inference for variables and literals, but it does
not try to infer function signatures or perform overload resolution.

The third design goal is a clear route to native code. The language has C-like
integers, raw pointers, and external declarations because these map naturally to
LLVM and to platform linkers. The compiler does not contain a full linker. It
emits an object file and expects a system C compiler to produce the final
executable.

The fourth design goal is compile-time execution. \epl{} should be able to
evaluate selected code during compilation without turning the language into a
separate macro system. A \texttt{comptime} expression is ordinary \epl{} code
with extra purity restrictions.

\section{Program Structure}

An \epl{} source file is a sequence of top-level items. The current item kinds
are function definitions, function declarations, and struct definitions.
Functions and types are stored in separate namespaces. Function declarations are
collected before function bodies are checked, so functions may call functions
that are declared later in the file and may be recursive. Struct definitions are
processed in source order, so a struct field may only refer to a type already
known at the point of the struct definition.

A function definition has a name, parameters, an optional return type, and a
body. A function declaration has the same signature form but ends with a
semicolon instead of a body. Function declarations are how \epl{} names external
functions. Listing \ref{lst:hello-source} shows the standard hello-world
example.

\begin{lstlisting}[language=EPL,caption={Hello world in EPL},label={lst:hello-source}]
fn main() -> i32 {
	printf("hello world\n");
	0
}

fn printf(fmt: *i8, ...);
\end{lstlisting}

The compiler enforces the entry point signature \texttt{fn main() -> i32}. The
language allows external variadic function declarations, as shown with
\texttt{printf}, but it does not allow defining variadic functions in \epl{}.
This is a deliberate implementation boundary: calling simple variadic C
functions is useful for examples, while implementing and lowering variadic
\epl{} function bodies would add ABI complexity.

\section{Types and Values}

Every expression in \epl{} has a statically known type after semantic analysis.
The built-in types are:

\begin{itemize}
  \item \texttt{unit}, the zero-sized type for expressions that produce no
        meaningful value;
  \item \texttt{bool}, with values \texttt{true} and \texttt{false};
  \item signed integers \texttt{i8}, \texttt{i32}, and \texttt{i64};
  \item unsigned integers \texttt{u8}, \texttt{u32}, and \texttt{u64};
  \item \texttt{ptr}, an opaque pointer type;
  \item \texttt{!}, the never type for expressions that do not complete
        normally.
\end{itemize}

Composite types are typed pointers, arrays, and user-defined structs. A typed
pointer is written \texttt{*T}. An array type is written \texttt{[T; N]}, where
\texttt{N} is currently a number literal. Structs are declared by name and store
their fields by value.

Integer literals may have suffixes such as \texttt{10u32} or \texttt{0i64}. If a
literal has no suffix and the surrounding context expects an integer type, the
literal uses that expected type. Otherwise, unsuffixed integer literals default
to \texttt{i32}. String literals have type \texttt{*i8} and are intended for
simple C-style external calls.

The literal \texttt{undefined} produces an undefined value of an expected type.
It therefore requires context. For example, \texttt{let x: i32 = undefined;} is
valid, while \texttt{let x = undefined;} is not, because the compiler cannot
infer the type from the value.

\section{Variables and Scope}

Local variables are introduced with \texttt{let}. Variables are mutable. A
\texttt{let} declaration can include an explicit type, an initializer, or both:

\begin{lstlisting}[language=EPL,caption={Variable declaration forms}]
let a = 10;
let b: u64 = 20;
let c: [i32; 4];
\end{lstlisting}

Blocks introduce lexical scopes. Name lookup starts in the current scope and
then searches outer scopes. A later declaration may shadow an earlier binding in
the same or an inner scope. Function arguments behave as mutable local variables
inside the function body because lowering turns each argument into an ordinary
local variable initialized from an argument expression. This choice simplifies
assignment to arguments and keeps all local storage in one model.

Reading a local value before assigning to it is undefined. The compiler reserves
storage for declarations without initializers, but it does not implement a
definite-assignment analysis. This is another example of a conscious scope
decision: the language can express low-level uninitialized storage, but safety
checks around it are not yet implemented.

\section{Expressions and Blocks}

\epl{} is expression-oriented. A block evaluates its statements in order. If it
ends with a final expression, the block's type is the type of that expression. If
there is no final expression, the block has type \texttt{unit}. Expression
statements discard their value.

The following function uses a block as the result of an \texttt{if} expression:

\begin{lstlisting}[language=EPL,caption={Expression-oriented conditional}]
fn abs_i32(x: i32) -> i32 {
	if x < 0 {
		-x
	} else {
		x
	}
}
\end{lstlisting}

The compiler checks that the branches of an \texttt{if} expression have
compatible types. An \texttt{if} without \texttt{else} is only accepted when a
\texttt{unit} result is expected. This prevents accidentally using a
non-exhaustive conditional where a value is required.

\section{Control Flow}

The language supports \texttt{return}, \texttt{break}, \texttt{continue},
\texttt{if}, \texttt{loop}, \texttt{while}, and \texttt{for}. The \texttt{return}
expression exits the current function and has type \texttt{!}. \texttt{break}
and \texttt{continue} are valid only inside loops and also have type
\texttt{!}. A \texttt{loop} expression can produce a value if it exits through
\texttt{break value}. A \texttt{while} expression has type \texttt{unit}.

\texttt{for} loops currently support half-open integer ranges of the form
\texttt{start..end}. The start and end expressions are evaluated once before the
loop. The compiler lowers a \texttt{for} loop into a block containing hidden
counter and target variables plus a normal \texttt{loop}. The visible loop
variable is assigned a per-iteration copy, so assigning to it does not modify the
hidden counter.

\begin{lstlisting}[language=EPL,caption={Loop forms used in EPL}]
fn sum_to(n: u32) -> u32 {
	let total: u32 = 0;
	for i in 0u32..n {
		total += i;
	}
	total
}
\end{lstlisting}

The lowering strategy is important because it keeps the rest of the compiler
small. After semantic analysis, there is no special lower IR operation for
\texttt{for}. There is only block, store, comparison, arithmetic, and loop.

\section{Operators and Casts}

Arithmetic operators are supported for integer operands only. Both operands must
have the same integer type. Supported operations are addition, subtraction,
multiplication, division, remainder, and unary negation for signed integers.
Compound assignment operators modify a place expression in place.

Comparison operators return \texttt{bool}. Integers support all comparison
operators. Booleans support equality and inequality. Pointer comparisons are not
currently implemented.

Logical \texttt{\&\&} and \texttt{||} require boolean operands and short-circuit.
They are lowered to conditional expressions. For example, \texttt{a || b}
returns \texttt{true} when \texttt{a} is true and otherwise evaluates
\texttt{b}. This is more explicit than carrying logical operators through every
compiler stage.

The supported casts are integer-to-integer and pointer-to-pointer. Narrowing
integer casts truncate. Widening integer casts sign-extend signed source values
and zero-extend unsigned source values. Pointer-to-pointer casts preserve the
pointer value. Pointer-to-integer and integer-to-pointer casts are not supported.

\section{Structs and Arrays}

Struct definitions list named fields. Field names must be unique within a
struct. Struct values are passed, returned, assigned, and stored by value.
Initializers must provide every field exactly once, although the initializer
field order does not need to match the declaration order.

\begin{lstlisting}[language=EPL,caption={Struct declaration and initialization}]
struct Vec2 {
	x: i32,
	y: i32,
}

fn main() -> i32 {
	let vec = Vec2 .{ x: 10, y: 20 };
	0
}
\end{lstlisting}

Arrays are fixed-size values. Array indexing requires a \texttt{u64} index.
Indexing does not perform bounds checks. Like structs, arrays are stored by
value. The compiler computes array layout from the element type layout and the
literal length in the type.

\section{Pointers and External Functions}

\epl{} exposes raw typed pointers. The expression \texttt{\&place} takes the
address of a place. The expression \texttt{ptr.*} dereferences a typed pointer.
The opaque \texttt{ptr} type cannot be dereferenced because it has no pointee
type. The language does not implement lifetime checking, borrow checking, null
checking, pointer arithmetic, alias analysis, or bounds checking.

Pointers are needed both as a systems programming feature and as a bridge to
external C functions. The examples use declarations such as:

\begin{lstlisting}[language=EPL,caption={External declarations}]
fn printf(fmt: *i8, ...);
fn exit(code: i32) -> !;
\end{lstlisting}

These declarations are emitted as LLVM function declarations. The final system
link step resolves the symbols. The C ABI support should be treated as limited:
simple integer and pointer calls work for the examples, but by-value aggregate
ABI details, exact C \texttt{void} behavior, platform-specific variadic
requirements, and complex signatures are not fully covered.

\section{Compile-Time Evaluation}

The \texttt{comptime} keyword asks the compiler to evaluate an expression during
compilation. The result is inserted into the typed tree as a constant before
lower IR generation. Listing \ref{lst:comptime-primes-short} shows a simplified
example derived from the project examples.

\begin{lstlisting}[language=EPL,caption={Computing values at compile time},label={lst:comptime-primes-short}]
let primes = comptime {
	let lut: [i32; 20];
	let ready = 0;
	let next = 3;
	while ready < 20 {
		while !is_prime(next) {
			next += 2;
		}
		lut[ready as u64] = next;
		ready += 1;
		next += 2;
	}
	lut
};
\end{lstlisting}

A compile-time expression may use literals except strings, arithmetic,
comparison, boolean operations, integer casts, blocks, conditionals, loops,
local variables, local mutation, arrays, structs, \texttt{break},
\texttt{continue}, and calls to \texttt{@pure} functions. It may not access
surrounding run-time variables, use strings, take addresses, dereference
pointers, call impure functions, or return from the enclosing function.

This design is the main language-level feature that distinguishes \epl{} from a
minimal C-like compiler. It allows non-trivial computation during compilation
without adding a separate macro language.
